% Part: lambda-calculus
% Chapter: lambda-definability
% Section: partial-recursive-functions

\documentclass[../../../include/open-logic-section]{subfiles}

\begin{document}

\olfileid{lam}{ldf}{par}
\olsection{Fungsi Rekursif Parsial iku \usetoken{S}{lambda definable}}

Fungsi rekursif parsial yaiku fungsi kang diolehake saka fungsi
dhasar kanthi komposisi, rekursi primitif, lan minimisasi tanpa
wates. Bedane saka fungsi rekursif umum yaiku fungsi kang digunakake
ing panelusuran tanpa wates ora kudu reguler. Tanpa syarat
regularitas, fungsi kang ditetepake kanthi minimisasi bisa wae
ora nduweni nilai kanggo sawetara argumen.

Sepisanan katon kaya-kaya cara kang padha kanggo nuduhake yen
kabeh fungsi rekursif umum (kang total) iku !!{lambda definable}
uga bisa digunakake kanggo mbuktekake yen kabeh fungsi rekursif
parsial iku !!{lambda definable}. Tuladhane, komposisi $f$ karo $g$
!!{lambda define}d dening $\lambd[x][F (G x)]$ yen $f$ lan $g$
!!{lambda define}d dening suku $F$ lan~$G$. Nanging kanggo fungsi
parsial, iki dadi masalah. Yen $g(x)$ ora nduweni nilai, tegese
$G x$ ora nduweni wujud normal. Biasane iki uga ndadekake $F (G
x)$ ora nduweni wujud normal, kaya kang dikarepake. Nanging
upamakna $F$ iku $\lambd[x][\lambd[y][y]]$; ing kene $F (G x)$
malah nduweni wujud normal ($\lambd[y][y]$).

Masalah iki isih bisa diatasi, lan ana cara kanggo
!!{lambda define} kabeh fungsi rekursif parsial kanthi nilai kang
ora katetepake diwakili dening suku tanpa wujud normal. Nanging
cara-cara iku rada luwih rumit lan kurang gampang dipahami
tinimbang pendekatan kang wis digunakake kanggo fungsi rekursif
umum. Ing kene kita mung nyathet teoremane tanpa bukti:

\begin{thm}
  Kabeh fungsi rekursif parsial iku !!{lambda definable}.
\end{thm}

\end{document}
